\chapter{Dataframes}\label{ch:pl:dataframes}

Two researchers compute five-minute bars from the same trades and get different volumes. One library names a bar by its left
edge and puts a trade stamped exactly on a boundary into the bar that starts there; the other names it by its right edge and
puts the same trade into the bar that ends there; one of them drops the minutes in which nothing traded, the other does not.
Both are right by their own conventions, and a signal that compares the two sets of bars finds a pattern that is nothing but the
convention. The \omterm{def:pl:dataframes:dataframe}{dataframe} is the research platform's working language; its idioms --- bars, as-of joins, rolling windows,
cross-sectional ranks, pivots --- are where its silent bugs live, and its engines differ enormously in what they cost. This chapter
writes the idioms once with explicit conventions for three engines (\texttt{firm.marketdf}), checks that the engines agree, and
runs one research pipeline over a month of quotes and trades eagerly, lazily and out of core.

\section{The dataframe model}

\begin{definition}[Dataframe]\label{def:pl:dataframes:dataframe}
A \emph{dataframe}\index{dataframe} is a table of named, typed columns of equal length, with operations that act on whole columns
and groups of rows: selection, filtering, computed columns, grouping and aggregation, joins, sorting, windows and reshaping.
\end{definition}

The three engines of the chapter implement the same model differently. pandas holds each column as a numpy array (or an Arrow
array) with a row index attached, and executes each operation when it is called. Polars holds columns in Arrow's memory layout,
has no row index, and can either execute each operation or build a plan and execute it later. DuckDB is a SQL engine: a query
is a plan, executed when its result is asked for (chapter~5). The \omterm{def:pl:columnar-formats:layouts}{columnar layout} of chapter~3 underlies all three.

\begin{definition}[Long format, wide format]\label{def:pl:dataframes:shape}
A table in \emph{long format}\index{long format} has one row per observation and a column for each variable, including the keys
(symbol, time, value). The same data in \emph{wide format}\index{wide format} have one row per value of one key and one column per
value of another (a row per minute, a column per symbol).
\end{definition}

\begin{omfigure}
\begin{tikzpicture}[font=\scriptsize,cell/.style={draw,minimum width=1.05cm,minimum height=0.42cm,inner sep=1pt}]
\foreach \r/\t/\s/\v in {0/09:31/AAA/10.1,1/09:31/BBB/20.4,2/09:32/AAA/10.2,3/09:32/BBB/20.3} {
  \node[cell] at (0,-\r*0.42) {\t}; \node[cell] at (1.05,-\r*0.42) {\s}; \node[cell] at (2.1,-\r*0.42) {\v};}
\foreach \x/\h in {0/time,1.05/symbol,2.1/close} {\node[cell,fill=omDef!15] at (\x,0.42) {\h};}
\node[font=\scriptsize\itshape] at (1.05,1.05) {long: a row per observation};
\draw[->,thick] (3.1,-0.4) -- node[above]{pivot} (4.6,-0.4);
\draw[<-,thick] (3.1,-0.9) -- node[below]{melt} (4.6,-0.9);
\foreach \r/\t/\a/\b in {0/09:31/10.1/20.4,1/09:32/10.2/20.3} {
  \node[cell] at (5.6,-\r*0.42) {\t}; \node[cell] at (6.65,-\r*0.42) {\a}; \node[cell] at (7.7,-\r*0.42) {\b};}
\foreach \x/\h in {5.6/time,6.65/AAA,7.7/BBB} {\node[cell,fill=omMeth!15] at (\x,0.42) {\h};}
\node[font=\scriptsize\itshape] at (6.65,1.05) {wide: a row per minute, a column per symbol};
\end{tikzpicture}

\omcaption{The same closing prices in \omterm{def:pl:dataframes:shape}{long format} (keys as columns) and in \omterm{def:pl:dataframes:shape}{wide format} (one key as rows, the other as columns).
Research code moves between the two: long for storage, joins and grouping; wide for cross-sectional arithmetic and
matrices.}\label{fig:pl:dataframes:shape}
\end{omfigure}

Market data are stored long (chapter~4's partitions are long tables), and most transformations are clearer long: a rolling window
is a window over the rows of one symbol, a cross-sectional rank a rank over the rows of one minute. The \omterm{def:pl:dataframes:shape}{wide format} serves the
matrix operations of portfolio construction (One Quant Book~7, chapters~24--26), and the pivot between the two is where missing
values appear: a minute in which one symbol did not trade is a missing cell in the wide table, never a row in the long one.

\section{Eager and lazy evaluation}

\begin{definition}[Eager evaluation, lazy evaluation]\label{def:pl:dataframes:eager}
An engine with \emph{eager evaluation}\index{eager evaluation} executes each operation when it is called and materialises its full
result. An engine with \emph{lazy evaluation}\index{lazy evaluation} records the operations as a plan and executes the plan only when
a result is requested, which lets it see the whole computation before running any of it.
\end{definition}

\omterm{def:pl:dataframes:eager}{Eager evaluation} is easy to reason about and to debug step by step; every intermediate table exists, and every intermediate table
costs memory. \omterm{def:pl:dataframes:eager}{Lazy evaluation} gives the engine the whole pipeline, which is what makes optimisation possible.

\section{Query plans and optimisation}

\begin{definition}[Query optimiser]\label{def:pl:dataframes:optimiser}
A \emph{query optimiser}\index{query optimiser} rewrites a plan into an equivalent cheaper one before executing it: it pushes filters
and projections toward the scans (so that fewer rows and columns are read, chapter~3), removes unused columns and operations, reorders
joins and aggregations, and chooses algorithms.
\end{definition}

A lazy query for one symbol's prices on one day, written as a scan of every trade file followed by a filter and a selection, is
optimised by Polars into a plan that reads one partition and four of six columns, and applies the filter while scanning
(\cref{fig:pl:dataframes:plan}).

\begin{omfigure}
\begin{tikzpicture}[font=\scriptsize,op/.style={draw,rounded corners=2pt,align=center,minimum width=4.4cm,minimum height=0.62cm,
  inner sep=2pt},>=Stealth]
\node[font=\small\bfseries] at (0,0.75) {as written};
\node[op,fill=omMeth!12] (a1) at (0,0) {select \texttt{ts}, \texttt{price}};
\node[op,fill=omMeth!12] (a2) at (0,-1.0) {filter symbol $=$ S007, date $=5$};
\node[op,fill=omThm!12] (a3) at (0,-2.0) {scan all 20 date partitions\\ all 6 columns};
\draw[->] (a3) -- (a2); \draw[->] (a2) -- (a1);
\draw[->,very thick,gray] (2.6,-1.0) -- node[above,black]{optimiser} (4.2,-1.0);
\node[font=\small\bfseries] at (6.8,0.75) {as executed};
\node[op,fill=omMeth!12] (b1) at (6.8,0) {select \texttt{ts}, \texttt{price} (2 of 2)};
\node[op,fill=omDef!12,minimum height=1.62cm] (b2) at (6.8,-1.5) {scan \texttt{date=5} only\\ project 4 of 6 columns\\
  filter symbol $=$ S007 while scanning};
\draw[->] (b2) -- (b1);
\node[anchor=north,align=center,text width=4.6cm] at (6.8,-2.45) {partition pruning, projection and predicate pushdown};
\end{tikzpicture}

\omcaption{A narrow query as written and as Polars executes it (from the plan it prints). The filter on the partition column removes 19 of
20 partitions before any file is opened; the projection keeps the two columns returned and the two the filter reads; the filter on the
symbol runs inside the scan.}\label{fig:pl:dataframes:plan}
\end{omfigure}

The partition pruning (only \texttt{date=5} is scanned), the \omterm{def:pl:columnar-formats:pushdown}{projection pushdown} (four columns: the two asked for and the two the filter
needs) and the \omterm{def:pl:columnar-formats:pushdown}{predicate pushdown} of chapter~3 are all the optimiser's work; an eager version of the same code would read the month and
filter it in memory.

\section{Out-of-core and streaming execution}

\begin{definition}[Out-of-core processing, streaming execution]\label{def:pl:dataframes:outofcore}
\emph{Out-of-core processing}\index{out-of-core processing} computes a result over data larger than memory by reading it in pieces
from storage. \emph{Streaming execution}\index{streaming execution} is the engine-level form of it: the plan's operators pass batches
to one another as they are produced, so that a scan, a filter and an aggregation hold one batch at a time, while operators that need
all of their input (a global sort, some joins) buffer it.
\end{definition}

Chapter~8 wrote out-of-core loops by hand; a lazy engine with a streaming executor does the same for a whole plan. It cannot do it
for every operator: an as-of join needs both sides sorted, and a sort needs all its input, so a plan containing them buffers at those
points. The pipeline below has both, which is why streaming saves memory here but does not make it constant.

\section{Idioms for market data}

\begin{definition}[Window function]\label{def:pl:dataframes:window}
A \emph{window function}\index{window function} computes, for each row, a value over a set of related rows --- the rows of the same
group, ordered, possibly limited to a frame such as the last $n$ rows --- without collapsing the rows as an aggregation does:
a rolling mean per symbol, a rank within a minute, the previous row's value (a lag).
\end{definition}

Five idioms cover most of research on market data, and each has a convention that must be stated rather than inherited.

\begin{method}[Market-data idioms with their conventions]\label{met:pl:dataframes:idioms}
\textbf{Bars}: state the width, the label (left or right edge), the closed side, and whether empty bars are dropped or filled (with
the previous close and zero volume). \textbf{As-of joins}: state the key (time or sequence number, chapter~4), the group, and whether
an equal time matches. \textbf{Rolling windows}: state the frame (rows or time), the group and the minimum number of observations.
\textbf{Cross-sectional ranks}: state the tie rule. \textbf{Pivots}: state what a missing cell means.
\end{method}

\omcode{code/firm/marketdf/firm_marketdf.py}{54}{88}{Bars with explicit conventions in three engines: the bar index from the closed
side by floor division (whatever the engine's integer division does with negative numbers), the label from the edge, then the same
aggregation.}\label{lst:pl:dataframes:bars}

\Cref{lst:pl:dataframes:bars} is the bars idiom in the three engines. Two details are the kind that breaks equivalence silently: the
first and last price of a bar must be taken in time order with a tie-breaker (the feed's sequence number), and the index of a
right-closed bar is $\lceil t/w \rceil - 1 = \lfloor (t-1)/w \rfloor$ for integer times, which each engine must compute with floor
division --- DuckDB's integer division truncates toward zero, and the SQL version computes the floor explicitly. The acceptance tests
compare the three engines on trades stamped to the whole second, so that boundaries and ties occur.

\begin{example}[Two sets of bars]\label{ex:pl:dataframes:bars}
On one symbol-day of synthetic trades (375 trades, stamped to the second as some vendors deliver them), five-minute bars labelled
and closed on the left and on the right differ in volume in 4 of the 74 intervals: two trades fall exactly on a five-minute boundary,
and each moves its volume, 600 shares in all, from one bar to the next. The two sets also carry labels five minutes apart for the same
interval. Filling empty bars adds 4 bars to the 74, and the day's bar returns become 77 instead of 73: a return computed across a
dropped bar spans ten minutes and is compared, in a rolling window, with five-minute returns.
\end{example}

\begin{omfigure}
\begin{tikzpicture}[font=\scriptsize,>=Stealth]
\draw[->] (0,0) -- (11.2,0) node[right]{time};
\foreach \x/\l in {0.5/10:00,5.5/10:05,10.5/10:10} {\draw (\x,0.1) -- (\x,-0.1) node[below]{\l};}
\fill[omThm] (5.5,0) circle (2.5pt);
\node[omThm,above] at (5.5,0.1) {trade at 10:05:00};
\draw[omDef,thick] (0.5,0.9) -- (5.4,0.9); \draw[omDef,line width=2pt] (5.6,0.9) -- (10.5,0.9);
\fill[omDef] (0.5,0.9) circle (2pt); \draw[omDef,fill=white] (5.4,0.9) circle (2pt);
\fill[omDef] (5.6,0.9) circle (2pt); \draw[omDef,fill=white] (10.5,0.9) circle (2pt);
\node[omDef,above] at (3,0.95) {bar labelled 10:00}; \node[omDef,above] at (8,0.95) {bar labelled 10:05 (holds the trade)};
\node[omDef,anchor=east,font=\scriptsize\bfseries] at (0.3,0.9) {left/left};
\draw[omMeth,line width=2pt] (0.5,-0.9) -- (5.4,-0.9); \draw[omMeth,thick] (5.6,-0.9) -- (10.5,-0.9);
\draw[omMeth,fill=white] (0.5,-0.9) circle (2pt); \fill[omMeth] (5.4,-0.9) circle (2pt);
\draw[omMeth,fill=white] (5.6,-0.9) circle (2pt); \fill[omMeth] (10.5,-0.9) circle (2pt);
\node[omMeth,below] at (3,-0.95) {bar labelled 10:05 (holds the trade)}; \node[omMeth,below] at (8,-0.95) {bar labelled 10:10};
\node[omMeth,anchor=east,font=\scriptsize\bfseries] at (0.3,-0.9) {right/right};
\end{tikzpicture}

\omcaption{A trade stamped exactly 10:05:00 under two bar conventions. Closed on the left, it belongs to the bar $[10\!:\!05,
10\!:\!10)$, labelled 10:05; closed on the right, to $(10\!:\!00, 10\!:\!05]$, also labelled 10:05 because that convention names bars by
their right edge. Filled dots mark closed ends, open dots open ends; the thick bar holds the trade. The same label names two different
intervals, and the trade's volume changes interval.}\label{fig:pl:dataframes:edge}
\end{omfigure}

\section{Joins, windows and ranks: where ties hide}

Bars are the most visible convention; three others hide in the joins and windows. An as-of join on the timestamp matches each trade
with the last quote at or before its time: when a trade and the quote change it caused share a timestamp, the inclusive join returns the
quote after the trade. Chapter~4 found this on 223 of 17\,281 trades and made the join strict on the feed's sequence number;
\texttt{firm\_marketdf.asof} takes \texttt{strict} as an argument in all three engines (pandas \texttt{allow\_exact\_matches}, Polars the
same, DuckDB the inequality \texttt{>} instead of \texttt{>=}).

A rolling window has a frame in rows or in time, and a minimum number of observations below which it returns nothing or a partial
value: the idiom's rolling mean averages whatever rows exist, one at the start of each symbol, which must be known to the code that uses
it. A cross-sectional rank has a tie rule, and engines differ in their default.

\begin{example}[Tied returns]\label{ex:pl:dataframes:ties}
In one minute, four symbols have returns $0.1$, $0.2$, $0.2$ and $0.3$ (in basis points). The average rule gives ranks 1,
2.5, 2.5 and 4, the rank sum $10 = 4\cdot5/2$ whatever the ties. SQL's \texttt{rank()} gives 1, 2, 2, 4 (the minimum rule), which lowers
the tied symbols' average rank; with $c$ rows tied at \texttt{rank()} $= r$, the average rank is $(2r + c - 1)/2 = (4 + 2 - 1)/2 = 2.5$,
the expression the SQL of \cref{lst:pl:dataframes:sql} uses. Returns of prices on a tick grid tie often: in the smaller month of the
tests, 10.8\% of the 222\,621 one-minute returns are exactly zero, and 6\,467 groups of tied returns occur in 7\,808 minutes.
\end{example}

\section{One pipeline, five ways}

The chapter's research pipeline runs over the chapter-3 month of quotes (now 3 million, three thousand per symbol per day) and a
month of 1.2 million trades, stored as date partitions of Parquet: one-minute bars, their log returns, the cross-sectional rank of
each return within its minute, each symbol's average rank over the month; and, in parallel, every trade joined as of its prevailing
quote and each symbol's mean effective spread in basis points. \Cref{lst:pl:dataframes:lazy} is the Polars version, a lazy plan
that the three Polars modes share.

\omcode{code/platforms/10-dataframes/python/pl_dataframes.py}{111}{124}{The pipeline as one Polars plan: bars, returns and ranks by
window expressions, and the as-of join and effective spread; the same plan is collected eagerly, lazily or by the streaming
engine.}\label{lst:pl:dataframes:lazy}

\omcode{code/platforms/10-dataframes/python/pl_dataframes.py}{140}{161}{The same pipeline in DuckDB's SQL, over the Parquet files:
bars by an ordered aggregate, returns by a lag, ranks by a window with the tie rule written out, and the as-of
join.}\label{lst:pl:dataframes:sql}

All five modes return the same fifty average ranks and fifty mean spreads to nine decimals. \Cref{fig:pl:dataframes:engines} measures
them, each in its own process with one thread, memory counted as the process's high-water mark minus that of a process that only imports
the libraries (155~MB). pandas, eager and with a row index, reads the whole month into memory and holds every intermediate table: it
needs 895~MB more than the imports and 1.3~s. Polars eager holds the same tables in Arrow form: 769~MB. Lazily, it needs 654~MB, and
with the streaming engine 614~MB; the as-of join and the sorts keep the saving modest. DuckDB, executing the SQL over the files with its
own streaming operators, needs 317~MB, a third of pandas.

\begin{omfigure}
\begin{tikzpicture}
\begin{axis}[width=12cm,height=5.6cm,ybar,bar width=10pt,ymin=0,ymax=1000,ylabel={extra memory (MB)},
  xtick={0,1,2,3,4},xticklabels from table={figdata/platforms/10-dataframes/measured_engines.csv}{engine},
  tick label style={font=\small},xticklabel style={text width=1.7cm,align=center},label style={font=\small},grid=major,grid style={gray!20},enlarge x limits=0.15,
  nodes near coords,point meta=explicit symbolic,nodes near coords style={font=\footnotesize},table/col sep=comma]
\addplot[fill=omDef!45,draw=omDef] table[x expr=\coordindex,y=extra_mb,meta=seconds]{figdata/platforms/10-dataframes/measured_engines.csv};
\end{axis}
\end{tikzpicture}

\omcaption{Memory beyond the imports (bars) and wall time in seconds (labels) of the same research pipeline over a month of quotes and
trades, in five engine modes, each run in its own process with one thread. The five give identical answers. Measured on a laptop (Intel
Core Ultra 7 155H) under WSL2, machine otherwise idle. Data: \texttt{bench\_dataframes.py}.}\label{fig:pl:dataframes:engines}
\end{omfigure}

\begin{remark}[Measuring memory of a child process]\label{rem:pl:dataframes:rss}
The first version of the measurement read the child's peak resident set from \texttt{getrusage}, and every engine seemed to need about
the same memory: on Linux, resource usage is preserved across \texttt{exec}, so the figure can include the memory of the process the child was started
from. The high-water mark in
\texttt{/proc/self/status} belongs to the process's own address space; it is what the chapter reports.
\end{remark}

\begin{dated}{2026-09}{Engine versions}\label{dat:pl:dataframes:versions}
The chapter's measurements and plans use pandas 2.3.3, Polars 1.44.2, DuckDB 1.5.5 and PyArrow 25.0.1, the versions pinned in the series'
requirements. Plans, defaults and memory use change between versions: the plan of \cref{fig:pl:dataframes:plan} is what this Polars
version prints, and Polars falls back to its in-memory engine for operations not yet implemented in streaming.
\end{dated}

\section{Tutorial: bars, conventions and five engines}\label{tut:pl:dataframes:tut}

\begin{tutorial}
\textbf{Goal.} Compute bars under stated conventions in three engines, see what the conventions change, and run one pipeline in five
modes with identical answers. \textbf{End state:} \cref{ex:pl:dataframes:bars}, \cref{fig:pl:dataframes:engines} and the optimised plan
of the chapter.
\begin{tutsteps}
\item \textbf{Idioms}: \texttt{firm\_marketdf.bars}, \texttt{asof}, \texttt{rolling\_mean}, \texttt{xs\_rank}, \texttt{to\_wide},
\texttt{to\_long}; the acceptance tests check that the three engines agree.
\item \textbf{Conventions}: \texttt{pl\_dataframes.conventions()} for the symbol-day of \cref{ex:pl:dataframes:bars}.
\item \textbf{Plan}: \texttt{plan\_text()} prints the optimised lazy plan of a narrow query.
\item \textbf{Pipeline}: \texttt{write\_month(**BIG)}, then \texttt{bench\_dataframes.py} runs the five modes, each in its own process,
checks the answers and records time and memory.
\end{tutsteps}
\textbf{What to change next.} Replace the as-of join in the Polars plan by a join on the previous second and see what the streaming
engine then saves; add a filled-bar option to the pipeline and measure how the average ranks change.
\end{tutorial}

\section{Build: market-data idioms}\label{bld:pl:dataframes:marketdf}

\begin{build}
\bfield{Purpose} The platform's shared implementation of the idioms every researcher rewrites: one contract, three engines, the
conventions as arguments, the equivalence as tests. The backtest engine of chapter~11 and the feature code of the research environment
(chapter~15) use it.
\bfield{Interface} \texttt{ENGINES}; \texttt{bars(ticks, width, label, closed, fill, engine)}; \texttt{asof(left, right, by, on,
engine, strict)}; \texttt{rolling\_mean(df, by, order, column, window, engine)}; \texttt{xs\_rank(df, time, column, engine)};
\texttt{to\_wide}, \texttt{to\_long}.
\bfield{Rules} Conventions are arguments with no hidden default beyond the documented one; every idiom has the same result in every
engine, to the row; ties are broken by the sequence number; integer times are bucketed by floor division in every engine.
\bfield{Acceptance tests} \texttt{code/firm/marketdf/tests/}: bars under four conventions equal across engines on second-stamped trades,
with volume conserved; the conventions worked by hand on four trades; as-of joins (inclusive and strict), rolling means and ranks equal
across engines; a pivot round trip.
\bfield{Stretch} Time-based rolling windows; bars with volume-weighted prices; the idioms for Arrow tables without pandas.
\end{build}

\begin{omsources}
\item H.~Wickham, ``Tidy data'', \emph{Journal of Statistical Software} 59(10), 2014.
\item pandas documentation, \texttt{DataFrame.resample}; Polars user guide, ``Streaming''; DuckDB documentation, ``AsOf join''.
\item Linux man-pages, \texttt{getrusage(2)}.
\end{omsources}

\section{Exercises}

\begin{exercise}[$\star$]\label{exo:pl:dataframes:1}
A trade is stamped 10:05:00 exactly. In which five-minute bar, and under which label, does it fall under each of the two conventions
of \cref{fig:pl:dataframes:edge}?
\end{exercise}

\begin{exercise}[$\star$]\label{exo:pl:dataframes:2}
Convert the long table of \cref{fig:pl:dataframes:shape} to wide and back. What does the wide table hold for a minute in which BBB did
not trade?
\end{exercise}

\begin{exercise}[$\star$]\label{exo:pl:dataframes:3}
Which columns and partitions does the optimised plan of the narrow query read, and why four columns rather than two?
\end{exercise}

\begin{exercise}[$\star\star$]\label{exo:pl:dataframes:4}
Why is the index of a right-closed bar $\lfloor (t-1)/w \rfloor$ for integer $t$, and what goes wrong in an engine whose integer division
truncates toward zero, for negative times?
\end{exercise}

\begin{exercise}[$\star\star$]\label{exo:pl:dataframes:5}
In \cref{ex:pl:dataframes:bars}, dropping the empty bars gives 73 returns and filling them 77. Which returns differ, and what does a
20-bar rolling volatility of the dropped series overstate?
\end{exercise}

\begin{exercise}[$\star\star$]\label{exo:pl:dataframes:6}
Why does \omterm{def:pl:dataframes:outofcore}{streaming execution} save less memory on the chapter's pipeline than on a pipeline of scans, filters and aggregations?
\end{exercise}

\begin{exercise}[$\star\star\star$]\label{exo:pl:dataframes:7}
\emph{Coding.} With \texttt{firm\_marketdf.bars}, compute one-minute bars of the chapter's month for one symbol in the three engines and
under the four conventions of the acceptance tests. Do the engines agree, and does the total volume depend on the convention?
\end{exercise}

\begin{exercise}[$\star\star\star$]\label{exo:pl:dataframes:8}
\emph{Find the flaw.} ``We resample trades to bars with the library default, compute returns with \texttt{pct\_change} on the bar table
sorted by time, and rank them across symbols each minute.''
\end{exercise}

\section{Problem: Two Sets of Bars}

\begin{problem}[{Weekend problem --- a month through five engines}]\label{pb:pl:dataframes:1}
The chapter's month of quotes and trades, its bar conventions and the five engine modes of \cref{fig:pl:dataframes:engines}.

\medskip\noindent\textbf{Part I --- Conventions.}
\begin{enumerate}
\item How many trades of the symbol-day fall exactly on a five-minute boundary, and why can that happen?
\item How many of the 74 intervals differ in volume between the two conventions, and how many shares move?
\item How far apart are the labels of the same interval under the two conventions?
\item How many bars and returns does the day have with empty bars dropped, and with them filled?
\item Which convention would you choose for a research store, and what must be written down?
\end{enumerate}

\medskip\noindent\textbf{Part II --- Plans.}
\begin{enumerate}[resume]
\item What does an eager engine materialise that a lazy one need not?
\item What three optimisations does the narrow query's plan show?
\item Why is a lazy plan easier to optimise than eager code?
\item Which operators of the pipeline cannot stream, and why?
\item Why must the SQL compute floor division explicitly?
\end{enumerate}

\medskip\noindent\textbf{Part III --- The five modes.}
\begin{enumerate}[resume]
\item What does the pipeline compute, and how are the five answers compared?
\item How much memory beyond the imports does each mode need?
\item Which mode is fastest, and which uses least memory?
\item Why is pandas the most expensive in memory?
\item Why did a first measurement show about the same memory for every engine?
\end{enumerate}

\medskip\noindent\textbf{Part IV --- The verdict.}
\begin{enumerate}[resume]
\item State the \emph{named result}: the volume and return discrepancies of the conventions, and the time and memory of the pipeline in
the five modes.
\item Which engine would you use for a month, and which for five years of the same data?
\item What would you standardise across a research team so that two researchers' bars agree?
\item What does an equivalence test across engines protect against?
\item In one sentence: what makes a \omterm{def:pl:dataframes:dataframe}{dataframe} pipeline trustworthy?
\end{enumerate}
\end{problem}

\section{Interview questions}

\begin{interviewq}[$\star$ \iqroles{researcher}]\label{iq:pl:dataframes:1}
Your five-minute bars do not match a colleague's. What do you check?
\end{interviewq}

\begin{interviewq}[$\star\star$ \iqroles{researcher, developer}]\label{iq:pl:dataframes:2}
What is the difference between eager and lazy \omterm{def:pl:dataframes:dataframe}{dataframe} evaluation, and when does it matter?
\end{interviewq}

\begin{interviewq}[$\star\star$ \iqroles{researcher}]\label{iq:pl:dataframes:3}
How would you compute a cross-sectional rank of one-minute returns across 3\,000 symbols for a year?
\end{interviewq}

\begin{interviewq}[$\star\star$ \iqroles{developer}]\label{iq:pl:dataframes:4}
Your pandas job runs out of memory on a month of trades. What do you change?
\end{interviewq}

\begin{interviewq}[$\star\star\star$ \iqroles{developer, researcher}]\label{iq:pl:dataframes:5}
Design a shared library of market-data transformations for a research team using several \omterm{def:pl:dataframes:dataframe}{dataframe} engines.
\end{interviewq}
